\section{Conclusion}
\label{sec:conclusion}

This paper develops an econophysics pipeline for analysing the dependency structure of Brazilian equities traded on B3 from 1998 to 2025. The main findings can be summarized along four dimensions.

First, the stylized-fact and correlation analysis provides descriptive evidence that Brazilian equities display empirical signatures commonly associated with complex financial systems: heavy-tailed return behaviour, volatility clustering, persistent absolute-return dependence and a moderate but positive average cross-correlation ($\bar{\rho}\approx0.24$). Within-sector correlations are larger than between-sector correlations on average, supporting the interpretation that sectoral organization is reflected in the empirical dependency matrix.

Second, Random Matrix Theory provides a spectral decomposition of the correlation matrix. Five eigenvalues exceed the Mar\v{c}enko--Pastur upper bound, with the largest eigenvalue accounting for approximately 37.3\% of the trace of the correlation matrix. The leading component is interpreted as a market-mode component, while the remaining eigenvectors outside the random band are consistent with commodity, financial and utility-related group structures. The reconstruction diagnostics indicate numerical accuracy within floating-point precision, and the group-mode matrix retains localized dependency patterns after the dominant market component is removed.

Third, financial network analysis through MST and PMFG representations shows that RMT filtering changes the topology of asset dependencies. The group-mode PMFG has lower mean edge correlation but higher clustering than the original PMFG, indicating stronger local clustering after market-mode removal. Hub rankings shift from networks dominated by broad market co-movement to networks where different assets occupy central positions after excluding the leading component. The PMFG retains 168 edges with triangles and 4-cliques, providing local topological structures that are absent from the MST and useful for mesoscopic analysis.

Fourth, the volatility forecasting experiment provides preliminary evidence that network topology can complement classical volatility predictors. GARCH(1,1) and HAR-RV remain strong benchmarks, especially at short horizons, consistent with the high persistence ($\alpha+\beta>0.98$) observed in all three target assets. Machine-learning models with PMFG-derived features approach these baselines, and at the 20-day horizon Ridge Regression obtains slightly lower average QLIKE than GARCH in this experiment. Since no formal equal-predictive-ability test is reported, this forecasting result should be interpreted descriptively. Feature importance analysis shows that classical autoregressive volatility features dominate, while selected PMFG centralities contribute smaller but nonzero predictive information.

The appropriate conclusion is incremental: network features do not replace classical volatility dynamics, but they add structural information from the cross-sectional geometry of the market. Future work should extend this framework by constructing rolling RMT and network features for real-time forecasting, expanding the target asset set, comparing against multivariate volatility models such as DCC-GARCH and evaluating temporal neural network and graph neural network architectures.
